\documentclass[11pt]{article}
\usepackage[margin=1in]{geometry}\usepackage[T1]{fontenc}
\usepackage{lmodern,amsmath,amssymb,amsthm,microtype}
\usepackage[colorlinks=true]{hyperref}
\newtheorem{theorem}{Theorem}[section]\newtheorem{lemma}[theorem]{Lemma}
\newcommand{\E}{\mathbb E}
\title{An exponential-depth sample bound for Gaussian Bayesian ReLU networks}
\author{ChatGPT-6 Astra \and Claude Opus 5.5 \and Kaarel H\"anni \and Dmitry Vaintrob\thanks{Authors are listed in alphabetical order.}}
\date{18 November 2026}
\hypersetup{pdftitle={An exponential-depth sample bound for Gaussian Bayesian ReLU networks},pdfauthor={ChatGPT-6 Astra; Claude Opus 5.5; Kaarel H\"anni; Dmitry Vaintrob}}
\begin{document}\maketitle
\begin{abstract}
We track depth in an agnostic approximation guarantee for a fixed Gaussian Bayesian neural network. The learner has hidden ReLU activations, a bounded final output used in both likelihood and prediction, and independent Gaussian fan-in priors with exponent at least two. At fixed teacher width, coefficient bound, input dimension, accuracy and confidence, $\exp(O(H))$ samples suffice at hidden depth $H$, uniformly over every learner width above the teacher width. A data-dependent Gaussian likelihood variance is chosen using two risk certificates. The proof normalizes quadratic energy by depth, retains the arithmetic--geometric mean bound on the product of layer norms, and keeps the Gaussian scale in the weighted-tail calculation. These steps avoid an unnecessary $H^H$ loss. A raw-readout extension gives the same exponential depth scale when $N^{\gamma-1}\ge H+1$. No stability restriction on prior support or prior over network sizes is imposed.
\end{abstract}
The purpose of this note is to identify the depth scale in a learning upper bound. Exponential depth dependence can arise from the capacity of unrestricted deep networks; a superexponential dependence can also be introduced accidentally by a norm estimate. The theorem below proves a genuinely exponential upper bound for a specified bounded-output Gaussian observation model. It does not assert polynomial dependence on depth, and its unbounded-readout extension has an additional width condition stated explicitly below.

The statistical ingredients are the smooth squared-loss Rademacher inequality of Srebro, Sridharan and Tewari~\cite{sst}, ordinary bounded-loss relative uniform convergence, and the ReLU parameter-count estimate~\cite{bartlett}. Every depth-dependent constant in the Gaussian comparison and calibration is tracked below. The two-certificate selection returns one posterior; it is not a mixture of learning rules.
\section{Exponential, rather than superexponential, dependence on depth}
\label{sec:normalized-depth}
This section tracks depth in the independent-Gaussian ReLU oracle. Fix input dimension $d$, teacher width $C$, coefficient cap $B\ge1$, accuracy $\varepsilon>0$, and confidence $\delta>0$. Only depth and learner width vary. Inputs lie in $[-1,1]^d$ and labels in $[-1,1]$; their joint law is otherwise arbitrary. A constant input coordinate is included, and internal biases, if desired, are counted through one additional constant wire. The comparison networks have at most $H$ hidden layers; identity padding makes their hidden depth exactly $H$ without increasing the coefficient cap.

Put $D=d+1$, $L=H+1$, and let the learner have hidden ReLU width $N\ge C$. Use independent centered Gaussian coordinates with variances
\[
 \operatorname{Var}(W^1_{ji})=N^{1-\gamma}/D,
 \quad \operatorname{Var}(W^\ell_{ji})=N^{-\gamma},
 \quad \operatorname{Var}(v_j)=N^{-\gamma-1},
 \qquad \gamma\ge2.
\]
The final output is $f_w=s(v^{\mathsf T}h^H)$, where $s$ clips to $[-1,1]$. It is used in both the Gaussian likelihood and prediction. Write $\omega=N^\gamma$ and
\[
 Q=\frac DN\|W^1\|_F^2+\sum_{\ell=2}^H\|W^\ell\|_F^2+N\|v\|_2^2,
 \qquad q=Q/L.
\]
For a positive parameter $\mu$, let
\[
 d\Pi_{S,\mu}=Z^{-1}\exp[-\omega L\widehat R_S(f_w)/(2\mu)]\,d\pi_N,
 \qquad F_S(\mu)=-\frac{2\mu}{\omega L}\log Z.
\]
Thus the ordinary Gaussian observation variance is $n\mu/(\omega L)$. This is a reparametrization of likelihood scale, not a change in prior.

\begin{theorem}[An all-width exponential-depth bound]
\label{thm:exp-depth-ReLU}
There is a data-dependent positive likelihood-variance rule, using only the fixed model, sample size and confidence, with the following property. For every fixed $d,C,B,\varepsilon,\delta$ there is a finite $a=a(d,C,B,\varepsilon,\delta)$, independent of $H,N,\gamma$, such that
\[
 n\ge\exp(a(H+1))
\]
implies, with probability at least $1-\delta$, simultaneously for all $N\ge C$ and $\gamma\ge2$,
\[
 R\left(\int f_w\,d\Pi_{S,\widehat\mu}\right)
 \le\inf_{g\in\mathcal G_{C,B,H}}R(s(g))+\varepsilon.
\]
Here $\mathcal G_{C,B,H}$ is the depth-at-most-$H$, width-at-most-$C$ ReLU comparison class with coefficients bounded by $B$ and the stated constant-wire convention. The learner chooses one posterior, not a mixture of widths or likelihoods.
\end{theorem}

The proof uses the same two-certificate method as the width-only oracle: an energy certificate at large widths and a parameter-count certificate at smaller widths. The new point is to normalize energy by depth before making tail and capacity estimates. Replacing a product of layer norms by $Q^{L/2}$, or discarding the Gaussian scale before taking a moment, would introduce an unnecessary $L^L$.

\subsection{Capacity retains the product of the layer energies}
Let $q_1=(D/N)\|W^1\|_F^2$, $q_\ell=\|W^\ell\|_F^2$ internally, and $q_L=N\|v\|_2^2$. Define the coordinate-first-layer path variation
\[
 P_1(w)=|v|^{\mathsf T}|W^H|\cdots|W^2|r^1,
 \qquad r^1_j=\sum_{k=0}^d|W^1_{jk}|.
\]
Cauchy--Schwarz at the first layer and spectral norm bounded by Frobenius norm internally give
\[
 |z_w(x)|\le P_1(w)\le\prod_{\ell=1}^L\sqrt{q_\ell}
             \le(Q/L)^{L/2}=q^{L/2}.
\]
The factors of $N$ and $D$ cancel exactly. ReLU homogeneity and exponential scalar contraction give the standard path-variation bound
\begin{equation}
 \widehat{\mathfrak R}_n(\{z_w:q(w)\le t\})
 \le\frac{k_H t^{L/2}}{\sqrt n},\qquad
 k_H=1+\sqrt{2\log(2D)}+\sqrt{2H\log2}.
 \label{eq:depth-capacity}
\end{equation}
For clarity, normalize each hidden feature by its forward absolute path mass. Each normalized feature lies in the iterated absolute convex hull of ReLU images of the coordinate atoms. Exponential contraction increases the moment-generating bound by at most a factor two at each layer. The initial coordinate maximum has Gaussian or Rademacher exponential bound with logarithm $\log(2D)$; optimizing the exponential parameter gives the displayed $k_H$. Multiplying by the final path mass proves the estimate. Saturation is a contraction, so the same bound applies to $f_w$.

\subsection{The normalized Gaussian tail has a depth-independent cutoff}
The exact prior law is
\[
 \omega Lq\sim\chi^2_{D_N},\qquad
 D_N=N(D+1)+(H-1)N^2,
 \qquad \frac{D_N}{\omega L}\le d_0:=D+1.
\]
The normalized Gaussian scale is at most $1/L$, rather than merely at most one. Set
\begin{equation}
 K_n=2d_0\log2+\frac4L\log(4n^4)+1,
 \qquad c_0=\frac18+\frac{d_0\log2}4.
 \label{eq:depth-tail-cutoff}
\end{equation}
For every sample and $\mu>0$, the posterior population-risk integral on
$\{\widehat R_S+\mu q>2F_S(\mu)+\mu K_n\}$ is at most $n^{-4}$.
Indeed, on $\widehat R_S+\mu q>u$,
\[
 e^{-\omega L\widehat R_S/(2\mu)}
 \le e^{-\omega Lu/(4\mu)}e^{\omega Lq/4}.
\]
The prior exponential moment is exactly $2^{D_N/2}$, and bounded population risk is at most four. After dividing by $Z=e^{-\omega LF_S/(2\mu)}$, the tail risk is bounded by
\[
 4\exp\left(\frac{\omega LF_S}{2\mu}-\frac{\omega Lu}{4\mu}
                                     +\frac{D_N\log2}2\right).
\]
Substitute $u=2F_S+\mu K_n$, use $D_N\le\omega Ld_0$ and $\omega\ge1$, and obtain $n^{-4}$. Also entropy duality gives
\begin{equation}
 \E_\rho q\le\frac4{\omega L}\operatorname{KL}(\rho\Vert\pi_N)+2d_0\log2.
 \label{eq:depth-entropy}
\end{equation}
Prior symmetry and bounded outputs give $F_S\le2$.

\subsection{An observable normalized-energy certificate}
Let $C_{\rm sm}$ be a universal constant in the smooth-loss Rademacher inequality. Put
\[
 J_n=\left\lceil\log_2(4n+K_n+2)\right\rceil,
 \quad t_n=\max\{1,\log(4(J_n+1)/\delta)\},
\]
\begin{equation}
 a_n=\max\left\{n^{-1/2},\,
 2^{L/2}C_{\rm sm}\left[
 \frac{\sqrt2k_H\log^{3/2}(2n)}{\sqrt n}
                   +2\sqrt{t_n/n}\right]\right\}.
 \label{eq:depth-statistical-constant}
\end{equation}
Increase $C_{\rm sm}$ by a fixed numerical factor if needed. Apply the smooth squared-loss bound to dyadic normalized-energy balls using \eqref{eq:depth-capacity}. On one event of probability at least $1-\delta/2$, simultaneously for all widths and exponents and all $q\le4n+K_n+1$,
\[
 R(f_w)-\widehat R_S(f_w)
 \le a_n\sqrt{\widehat R_S(f_w)}(1+q)^{L/2}
                            +a_n^2(1+q)^L.
\]
The factor $2^{L/2}$ pays for rounding a radius upward to its containing dyadic shell. There is no factor $L^{L/2}$.

Select the unique positive solution of
\begin{equation}
 \mu^L=256a_n^2[2F_S(\mu)+(K_n+1)\mu]^{L-1}.
 \label{eq:depth-root}
\end{equation}
Continuity and monotonicity of $F_S(\mu)/\mu$ prove existence, uniqueness and measurability. The root is at least $1/n$. Write $U=2F_S+\mu(K_n+1)$. Off the weighted-objective tail, $\widehat R_S\le U$ and $1+q\le U/\mu$. The two statistical error terms are at most $\mu(1+q)/16$ and $\mu(1+q)/256$, respectively. Integrating and using \eqref{eq:depth-entropy} absorbs their energy cost in the relative-entropy term of $F_S$. Thus
\begin{equation}
 \E_{\Pi_{S,\mu}}R(f_w)\le F_S(\mu)+c_0\mu+n^{-4}
                         =:B_{\rm imp}(S).
 \label{eq:depth-certificate}
\end{equation}
Jensen gives the same certificate for the ordinary posterior mean.

If $F_S(\mu)\le u+\mu T$ for all $\mu\ge1/n$, set
$\mathcal B=2T+K_n+1+c_0$. Splitting the root according to whether $\mu(2T+K_n+1)$ is at least $2u$ gives the explicit bound
\begin{equation}
 B_{\rm imp}\le u+
 4\mathcal B(256a_n^2)^{1/L}u^{1-1/L}
       +256a_n^2(2\mathcal B)^L+n^{-4}.
 \label{eq:depth-empirical-oracle}
\end{equation}
All numerical constants here are independent of depth.

\subsection{The exact selective comparison cost}
Let the teacher width be $C$ and its coefficient cap $B\ge1$. Positive-homogeneous layer balancing gives an implementation, padded to any width $N\ge C$, whose individual weighted layer energies are at most
\[
 e=D^2C^2B^2.
\]
To verify this, initially multiply all hidden states by $\sqrt N$ and divide the readout by $\sqrt N$. The product of the $L$ weighted layer energies is at most $D^2C^{2H}B^{2L}$. Equalizing nonzero layer energies by successive positive rescalings proves the claim; the zero-function case is immediate.

Use the selective Gaussian comparison that narrows only coordinates entering the $C$ active rows and the readout. Its narrowed coordinate count is
\[
 D_{\rm nar}=CD+(H-1)CN+N.
\]
With standard deviations multiplied by $r\le1$, the exact entropy bound is
\[
 \frac2{\omega L}\operatorname{KL}(\nu_r\Vert\pi_N)
 \le e+\frac{2D_{\rm nar}}{\omega L}\log(1/r).
\]
Keeping the individual layer bound $e$ in the perturbation recurrence gives output RMS error at most
\[
 rK_e,\qquad K_e=\sqrt2(H+1)(e+1)^{H/2}.
\]
For example, normalized hidden second moments are at most $2(e+1)^\ell$, and active-row RMS errors obey
$\Delta_\ell\le\sqrt e\Delta_{\ell-1}+\sqrt2r(e+1)^{(\ell-1)/2}$.
This is the same Gaussian comparison recursion as usual, with the per-layer energy retained rather than replaced by total energy $Le$.

Choose $r=(8nK_e)^{-1}$. Bounded-output squared loss is 4-Lipschitz in prediction, so
\begin{equation}
 F_S(\mu)\le\widehat R_S(s(g))+1/(2n)+\mu T_N,
 \quad
 T_N=e+\frac{2D_{\rm nar}}{\omega L}\log(8nK_e).
 \label{eq:depth-teacher-cost}
\end{equation}
If $N^{\gamma-1}\ge L$, then
\[
 \frac{D_{\rm nar}}{\omega L}
 \le\frac{C+(D+1)/L}{N^{\gamma-1}}
 \le\frac{C+D+1}{L}.
\]
Consequently, with $t=(\log n)/L$,
\begin{equation}
 T_N\le c_{d,C,B}(1+t),\qquad
 K_n\le c_d(1+t),\qquad
 \mathcal B\le c_{d,C,B}(1+t).
 \label{eq:depth-linear-log}
\end{equation}
These constants do not depend on depth or learner width.

\subsection{Finishing the all-width bound}
For $n=e^{tL}$ with $t\ge1$, \eqref{eq:depth-statistical-constant} gives
\[
 a_n^2\le c_{d,\delta}\,2^L(1+t)^4L^5e^{-tL}.
\]
The harmless fifth power allows all confidence and peeling logarithms. Since $L^{5/L}$ is uniformly bounded for $L\ge2$, \eqref{eq:depth-empirical-oracle} and \eqref{eq:depth-linear-log} show that its two excess terms tend to zero as $t\to\infty$, uniformly over all $L\ge2$. For instance the first is bounded by a fixed polynomial in $1+t$ times $e^{-t}$, and the second by
$c_{d,\delta}(1+t)^4L^5[c_{d,C,B}(1+t)e^{-t}]^L$.
A sufficiently large fixed $t\ge a(d,C,B,\varepsilon,\delta)$ therefore makes the excess less than $\varepsilon/3$. Bounded-loss concentration for one deterministic near-minimizing teacher adds less than $\varepsilon/3$, uniformly in $L$, after increasing $a$. This proves the desired bound whenever $N^{\gamma-1}\ge L$.

If $N^{\gamma-1}<L$, then $N<L$ because $\gamma\ge2$. The parameter count is at most $c_dL^3$ and its ReLU pseudodimension is at most $c_dL^4\log(eL)$. The positive-temperature parameter certificate uses the first dyadic $\mu$ for which $F_S(\mu)$ is within $1/n$ of the empirical loss infimum. Full support guarantees such a positive value. Relative uniform convergence gives an excess bounded by
\[
 c_d\sqrt{\frac{L^4\log(eL)\log(en)+\log((N+1)^2/\delta)}n}
 +c_d\frac{L^4\log(eL)\log(en)+\log((N+1)^2/\delta)}n.
\]
The prescribed teacher fits at every $N\ge C$. The sample budget $e^{aL}$ dominates this polynomial for a sufficiently large constant $a$. A summable width-confidence allocation gives a common event, and the logarithm involving $N$ is at most $\log L$ in this branch.

Finally choose the smaller observable certificate. The rule does not use $C$, $B$, the desired accuracy, or the population risk. The analysis alone splits the widths at $N^{\gamma-1}=L$. Allocating confidence among the common energy event, parameter event, and fixed-teacher concentration completes the proof of Theorem~\ref{thm:exp-depth-ReLU}.

\section{Input dimension proportional to depth}
The same exponential-depth conclusion holds when input dimension grows linearly with depth. This is useful when comparing depth upper and lower bounds whose input dimensions are themselves proportional to depth.

\begin{theorem}[A uniform bound for $d=O(H)$]
\label{thm:linear-dimension-depth}
Fix $\kappa,C,B,\varepsilon,\delta$. There is a finite constant
$a=a(\kappa,C,B,\varepsilon,\delta)$ such that the conclusion of Theorem~\ref{thm:exp-depth-ReLU} holds for every $H\ge1$ and every
$d\le\kappa(H+1)$ whenever $n\ge e^{a(H+1)}$.
The normalized-energy certificate uses the sharper constants specified below. No other learner assumption changes.
\end{theorem}
\begin{proof}
Two coarse constants in the preceding proof can be kept exact. First,
\[
 \frac{D_N}{\omega L}\le\frac{D+H}{L}
       =1+\frac dL\le1+\kappa,
\]
so use $d_0=1+d/L$ in the tail and entropy bounds.
Second, the product calculation in the teacher-balancing step gives the per-layer bound
\[
 e=D^{2/L}C^{2H/L}B^2.
\]
This is bounded by a constant depending only on $\kappa,C,B$ when
$D=d+1\le(\kappa+1)L$: the quantity $L^{2/L}$ is uniformly bounded for $L\ge2$. If some layer energy vanishes, the represented function is zero and that case requires no balancing.

The large-width comparison ratio is still
\[
 \frac{D_{\rm nar}}{\omega L}
 \le\frac{C+(D+1)/L}{N^{\gamma-1}},
\]
whose numerator is bounded in terms of $C,\kappa$. Therefore
$T_N,K_n,\mathcal B$ remain $O_{\kappa,C,B}(1+\log n/L)$ when
$N^{\gamma-1}\ge L$.
The capacity coefficient satisfies $k_H^2=O_\kappa(L)$ because
$\log D=O_\kappa(\log L)$. Hence the same bound
$a_n^2\le c_{\kappa,\delta}2^L(1+t)^4L^5e^{-tL}$ holds.
Every step making the root excess small uniformly in $L$ is unchanged.

In the complementary branch $N<L$, both the input dimension and width are $O_\kappa(L)$. Parameter count is $O_\kappa(L^3)$ and pseudodimension is $O_\kappa(L^4\log(eL))$, as before. The exponential sample budget dominates this polynomial. The same certificate selection and confidence allocation complete the proof.
\end{proof}

\section{Unbounded readouts at sufficiently large widths}
\label{sec:raw-large-width-depth}
The normalized-energy argument also gives an exponential-depth upper bound for a Gaussian likelihood on the raw ReLU readout. The only missing part of the preceding all-width proof is its small-width parameter certificate, which used bounded outputs. We therefore state the raw-output extension with an explicit additional width condition.

\begin{theorem}[Raw ReLU readouts above a depth-dependent width]
\label{thm:raw-large-width-depth}
Fix $\kappa,C,B,\varepsilon,\delta$. Use the same independent Gaussian prior and input convention, but let the likelihood and prediction use the raw readout $z_w$, with no final saturation. There is a finite constant $a$, depending only on those fixed quantities, such that for $d\le\kappa(H+1)$ and $n\ge e^{a(H+1)}$, the normalized-energy root below gives, with probability at least $1-\delta$,
\[
 R\left(\int z_w\,d\Pi_{S,\widehat\mu}\right)
 \le\inf_{g\in\mathcal G_{C,B,H}}R(g)+\varepsilon,
\]
simultaneously for all $\gamma\ge2$ and widths satisfying
\[
 N\ge C,\qquad N^{\gamma-1}\ge H+1.
\]
The Gaussian observation variance is $n\widehat\mu/[N^\gamma(H+1)]$. One ordinary raw posterior mean is used.
\end{theorem}
\begin{proof}
Retain $L=H+1$, $q=Q/L$, and $d_0=1+d/L\le1+\kappa$.
The capacity and amplitude bound are
\[
 \widehat{\mathfrak R}_n(\{z_w:q\le t\})\le k_Ht^{L/2}/\sqrt n,
 \qquad |z_w|\le q^{L/2}.
\]
Under the centered Gaussian prior, each ReLU layer halves the corresponding centered Gaussian second moment. Since $N/N^\gamma\le1$, the normalized hidden second moments are at most $2^{-\ell}$ and the output second moment is at most one. Readout symmetry therefore gives prior squared risk at most two, so $F_S\le2$ as before.

The raw loss obeys $R(z_w)\le2(1+q^L)$. Exponential tilting of the exact chi-square law gives
\begin{align*}
 \E_\pi[(1+q^L)e^{\omega Lq/4}]
 &=2^{D_N/2}\left[1+
       \left(\frac2{\omega L}\right)^L
                         \prod_{j=0}^{L-1}(D_N+2j)\right]\\
 &\le2^{D_N/2}M_L,
 \qquad M_L=1+(2d_0+4)^L.
\end{align*}
The last bound retains $1/(\omega L)$ rather than replacing it by one. Define
\[
 K_n^{\rm raw}=2d_0\log2+\frac4L\log(2M_Ln^4)+1.
\]
The weighted-objective proof now bounds raw posterior tail risk by $n^{-4}$. Its cutoff remains
$O_\kappa(1+\log n/L)$ because $\log M_L=O_\kappa(L)$.

On normalized-energy shells the raw squared-loss envelope is bounded by $4(1+q)^L$. The same smooth-loss argument therefore gives the statistical inequality used earlier, with the same form of $a_n$ after enlarging its numerical constant. Use $K_n^{\rm raw}$ in the peeling cutoff and in the unique positive root
\[
 \mu^L=256a_n^2[2F_S(\mu)+(K_n^{\rm raw}+1)\mu]^{L-1}.
\]
The entropy calculation is unchanged, so its raw posterior-average risk is bounded by $F_S+c_0\mu+n^{-4}$ and obeys the empirical oracle \eqref{eq:depth-empirical-oracle}.

The balanced teacher has per-layer energy
$e=D^{2/L}C^{2H/L}B^2=O_{\kappa,C,B}(1)$ and amplitude at most
$U=e^{L/2}$. Its selective comparison has output RMS error at most $rK_e$, with $K_e=\sqrt2(H+1)(e+1)^{H/2}$. Choose
\[
 r=[16n(U+1)K_e]^{-1}.
\]
The comparison's raw empirical-risk increment is at most
$2(U+1)rK_e+(rK_e)^2\le1/(2n)$.
Moreover $\log(1/r)=\log n+O_{\kappa,C,B}(L)$. The exact entropy count therefore still gives
$T_N=O_{\kappa,C,B}(1+\log n/L)$ whenever $N^{\gamma-1}\ge L$.

The comparison class includes the zero function, so choose one deterministic near-minimizer with population risk at most $1+\varepsilon/10$. Its bounded range is at most $(U+1)^2\le4(e+1)^L$. Scalar Bernstein concentration shows that its empirical risk differs from its population risk by at most $\varepsilon/10$ once $n\ge e^{aL}$ for a sufficiently large constant $a$. In particular the empirical comparator term is bounded by a constant independent of $L$. The two terms in the empirical root oracle now vanish uniformly in $L$ as $\log n/L$ grows, exactly as in the bounded-output proof. Increasing $a$ completes the stated risk bound.

Every raw posterior moment used here exists: the prior is Gaussian at fixed finite architecture, the ReLU output grows polynomially in its parameters, and the likelihood is bounded by one with a positive normalizer. Jensen therefore applies to the ordinary raw posterior mean. No small-width parameter argument is invoked.
\end{proof}

The additional condition is modest in the stronger Gaussian regimes: at $\gamma=2$ it is $N\ge H+1$, and for larger fixed exponents it is $N\ge(H+1)^{1/(\gamma-1)}$. It remains separate from the all-width bounded-output theorem; that theorem does not supply a raw-output conclusion below this threshold.

\begin{thebibliography}{9}
\bibitem{sst} Nathan Srebro, Karthik Sridharan and Ambuj Tewari. Smoothness, Low Noise and Fast Rates. \emph{Advances in Neural Information Processing Systems} 23, 2010.
\bibitem{bartlett} Peter L. Bartlett, Nick Harvey, Christopher Liaw and Abbas Mehrabian. Nearly-tight VC-dimension and pseudodimension bounds for piecewise linear neural networks. \emph{Journal of Machine Learning Research} 20(63):1--17, 2019.
\end{thebibliography}
\end{document}
